\slajdid{09_2-s049}
\begin{frame}[label=09_2-s049]{ESD zaštita}
\gusttekst\setlength{\abovedisplayskip}{3pt}\setlength{\belowdisplayskip}{3pt}\setlength{\abovedisplayshortskip}{2pt}\setlength{\belowdisplayshortskip}{2pt}\begin{columns}[c,onlytextwidth]\column{.32\textwidth}\centering\begin{tikzpicture}[x=.75cm,y=.8cm,font=\sitneoznake,>=Latex]
\ctikzset{bipoles/length=.65cm}
\draw(0,.25)--(0,.75);\draw(-.13,.22)--(-.13,.78);\draw(0,.7)--(.4,.7)--(.4,1.1);\draw(0,.3)--(.4,.3)--(.4,-.35)node[ground]{};
\draw(-.75,.5)--(-.13,.5);\draw(-.75,1.15)--(-2.15,1.15);\draw(-2.15,1.15)--(-2.33,1.28)--(-2.65,1.28)--(-2.83,1.15)--(-2.65,1.02)--(-2.33,1.02)--cycle;\node[left]at(-2.83,1.15){$V_I$};
\draw(-1.6,-.35)node[ground]{} to[D](-1.6,1.15) to[D](-1.6,2.65);\draw(-1.85,2.65)--(-1.35,2.65)node[above left]{$V_{DD}$};\node[right]at(.55,.5){N};
\draw(0,1.65)--(0,2.15);\draw(-.13,1.62)--(-.13,2.18);\draw(0,2.1)--(.4,2.1)--(.4,2.65);\draw(0,1.7)--(.4,1.7)--(.4,1.1);
\draw(-.22,1.9)circle(.09);\draw(-.31,1.9)--(-.75,1.9)--(-.75,.5);\node[right]at(.55,1.9){P};
\draw(.15,2.65)--(.65,2.65)node[above left]{$V_{DD}$};\draw(.4,1.15)--(1.1,1.15)node[ocirc]{}node[right]{$V_O$};
\end{tikzpicture}


\column{.66\textwidth}
Osnova rada CMOS logičkih kola jeste tanak oksid na gejtu tranzistora. Sam ulazak u gejt je visokoimpedansna tačka koja ako služi kao ulaz spoljnih signala nema putanju ni prema napajanju ni prema masi. Bilo koja relativno mala količina naelektrisanja dovedena na taj ulaz može izazvati velika polja na gejtu a time i probijanje oksida gejta. Zbog toga se na gejtove odnosno ulazne priključke obavezno stavljaju zaštitne diode čija je uloga da obezbede odvođenje viška naelektrisanja prema napajanju odnosno prema masi i da ne dozvole potencijale na gejtovima tranzistora većim od napona napajanja odnosno nižim od napona mase.\par\medskip
CMOS logička kola su izuzetno osetljiva na statičko naelektrisanje pa je jedno od opštih pravila da pinovi CMOS logičkih kola ne smeju da se dodiruju rukama. Takođe rukovanje sa CMOS logičkim kolima podrazumeva da i oprema i operater budu „uzemljeni“ odnosno na potencijalu mase, na kojem se nalazi i logičko kolo.
\end{columns}
\end{frame}
